\documentclass[11pt,a4paper]{article}

\usepackage{amsmath,amssymb,amsthm}
\usepackage[utf8]{inputenc}
\usepackage{mathrsfs}
\usepackage[margin=1in]{geometry}
\usepackage{hyperref}
\usepackage{enumitem}
\usepackage{booktabs}

% Theorem environments
\theoremstyle{plain}
\newtheorem{theorem}{Theorem}[section]
\newtheorem{proposition}[theorem]{Proposition}
\newtheorem{lemma}[theorem]{Lemma}
\newtheorem{corollary}[theorem]{Corollary}

\theoremstyle{definition}
\newtheorem{definition}[theorem]{Definition}
\newtheorem{example}[theorem]{Example}

\theoremstyle{remark}
\newtheorem{remark}[theorem]{Remark}

% Operators and shortcuts
\newcommand{\E}{\mathbb{E}}
\newcommand{\R}{\mathbb{R}}
\newcommand{\Z}{\mathbb{Z}}
\newcommand{\C}{\mathbb{C}}
\newcommand{\T}{\mathbb{T}}
\newcommand{\Ent}{\mathrm{Ent}}
\newcommand{\Var}{\mathrm{Var}}
\newcommand{\tr}{\mathrm{tr}}
\newcommand{\Hess}{\mathrm{Hess}}

\title{Bakry--\'Emery Curvature of the Fractional Laplacian\\
via Fractional Brownian Covariance}

\author{Ramiro Fontes\\
Quijotic Research\\
\texttt{ramirofontes@gmail.com}}

\date{February 2026}

\begin{document}

\maketitle

\begin{abstract}
We prove that the carr\'e du champ of the symmetric $\gamma$-stable
generator $L_\gamma = -(-\Delta)^{\gamma/2}$, evaluated on Fourier modes,
coincides with the covariance kernel of fractional Brownian motion (fBM)
with Hurst parameter $H = \gamma/2$:
$\Psi_\gamma(\xi,\eta) = R_{\gamma/2}(|\xi|,|\eta|)$.
Both formulas are classical; their identification appears to be new.

We exploit this correspondence to develop the first positive Bakry--\'Emery
curvature theory for the fractional Laplacian.
On the torus $\T = \R/(2\pi\Z)$, the iterated carr\'e du champ
kernel $\Gamma_{\gamma,2}$ is the Hadamard square of $\Gamma_\gamma$:
the curvature becomes the generalized eigenvalue problem
$R_H^{\circ 2}\,v = \kappa\,R_H\,v$ for fBM covariance matrices.
At $\gamma = 1$ (Cauchy process, $H = 1/2$), the eigenvalues are the
odd integers $\{1,3,\ldots,2N{-}1\}$, giving curvature $\kappa = 1$
exactly.
On real trigonometric polynomials, the cross-sign kernel
$\Psi_\gamma(n,-m)$ vanishes if and only if $\gamma = 1$---a
consequence of the linearity of $|\cdot|^1$---making the Cauchy
process the unique stability index with $\kappa \geq 1$.

For the L\'evy--Fokker--Planck equation with confining potential
$V(x) = -\omega^2\cos x$, we derive an exact curvature formula
whose parameters are controlled by the fBM covariance.
These results complement the negative result of Spener--Weber--Zacher
that $(-\Delta)^{\gamma/2}$ on $\R^d$ fails
$\mathrm{CD}(\kappa,N)$ for all finite~$N$.
\end{abstract}

\medskip
\noindent\textbf{MSC 2020:} 60G52, 60G22, 35S10, 47D07, 42A16.\\
\textbf{Keywords:} Bakry--\'Emery curvature; fractional Laplacian;
carr\'e du champ; fractional Brownian motion; Hadamard product;
stable semigroup; cross-sign vanishing; L\'evy--Fokker--Planck;
Schoenberg kernel.


%======================================================================
\section{Introduction}\label{sec:intro}
%======================================================================

\subsection{Motivation}

The Bakry--\'Emery $\Gamma_2$-criterion is the principal tool for
establishing functional inequalities for Markov
semigroups~\cite{BakryGentilLedoux2014}. For diffusion generators
$L = \Delta + \nabla V \cdot \nabla$ on a Riemannian manifold,
the condition $\Gamma_2(f,f) \geq \kappa\,\Gamma(f,f)$ is equivalent
to a Ricci curvature lower bound, and the constant $\kappa$ controls
convergence to equilibrium through
\[
  \Gamma_2 \geq \kappa\,\Gamma
  \;\Longrightarrow\;
  \text{log-Sobolev with } C = 2/\kappa
  \;\Longrightarrow\;
  \text{Poincar\'e with } \lambda_1 \geq \kappa.
\]
For non-local generators such as the fractional Laplacian
$L_\gamma = -(-\Delta)^{\gamma/2}$, the theory is substantially
more delicate. Spener, Weber, and
Zacher~\cite{SpenerWeberZacher2020} proved that the fractional
Laplacian on $\R^d$ fails to satisfy the curvature-dimension inequality
$\mathrm{CD}(\kappa,N)$ for all $\kappa \in \R$ and all finite $N > 0$.
This was identified as a key open problem in extending the $\Gamma$-calculus
to non-local operators~\cite{Garofalo2019}.

In this paper, we work on the compact torus $\T = \R/(2\pi\Z)$, where
the spectral gap $\lambda_1 = 1$ provides the essential compactness that
the $\R^d$ setting lacks.  Our starting point is the observation that
the Fourier-space carr\'e du champ of the $\gamma$-stable generator
coincides with the covariance kernel of fractional Brownian motion
with Hurst parameter $H = \gamma/2$.
This identification---a connection between the Dirichlet form of a
jump process and the covariance of a Gaussian process---allows us to
express the full Bakry--\'Emery curvature theory in terms of spectral
properties of fBM covariance matrices.


\subsection{The stable--fBM correspondence}

The Fourier-space kernel of the carr\'e du champ
$\Psi_\gamma(\xi,\eta) :=
\frac{1}{2}(|\xi|^\gamma + |\eta|^\gamma - |\xi-\eta|^\gamma)$
is a standard object in the theory of jump-type Dirichlet
forms~\cite{BoettcherSchillingWang2013}.
The fBM covariance $R_H(s,t) = \frac{1}{2}(|s|^{2H} + |t|^{2H}
- |s-t|^{2H})$ is classical since
Mandelbrot--Van~Ness~\cite{Mishura2008}.
The kernel $\frac{1}{2}(|\xi|^\gamma + |\eta|^\gamma
- |\xi-\eta|^\gamma)$ also appears in Schoenberg's
theory of conditionally negative definite
functions~\cite{Schoenberg1938}.

These three appearances have not, to our knowledge, been connected
explicitly.  We prove (Theorem~\ref{thm:duality},
Section~\ref{sec:duality}):
\begin{equation}\label{eq:duality-intro}
  \Psi_\gamma(\xi,\eta) = R_{\gamma/2}(|\xi|,|\eta|),
  \qquad \gamma \in (0,2),
\end{equation}
identifying the carr\'e du champ of the stable generator with
the covariance of fBM.  The proof requires only the
L\'evy--Khintchine formula.
The identity was originally discovered within an operator
factorization framework for stochastic
calculus~\cite{Fontes2026a}; here we give a self-contained
treatment and develop its consequences for curvature theory.


\subsection{Main results}

We work on $\T = \R/(2\pi\Z)$ with normalized Lebesgue measure
$\mu = dx/(2\pi)$. The fractional Laplacian acts on Fourier modes as
$L_\gamma(e^{inx}) = -|n|^\gamma\,e^{inx}$.

\medskip\noindent
\textbf{Theorem~A} (Hadamard square identity; \S\ref{sec:gamma2}).
\emph{The iterated carr\'e du champ kernel of $L_\gamma$ is the
Hadamard square of the carr\'e du champ kernel:
\begin{equation}\label{eq:intro-A}
  \Phi_{\gamma,2}(\xi,\eta)
  = \bigl[\Psi_\gamma(\xi,\eta)\bigr]^2
  = \bigl[R_{\gamma/2}(|\xi|,|\eta|)\bigr]^2.
\end{equation}}

\medskip\noindent
\textbf{Theorem~B} (Curvature as eigenvalue problem; \S\ref{sec:curv-pos}).
\emph{On positive-frequency polynomials $f = \sum_{n=1}^N a_n\,e^{inx}$,
the Bakry--\'Emery curvature is
\begin{equation}\label{eq:intro-B}
  \kappa(\gamma, N) = \lambda_{\min}\!\bigl(R_H^{\circ 2},\, R_H\bigr),
  \quad H = \gamma/2,
\end{equation}
where $R_H$ is the $N \times N$ fBM covariance matrix on $\{1,\ldots,N\}$.}

\medskip\noindent
\textbf{Theorem~C} (Odd integer eigenvalues; \S\ref{sec:cauchy}).
\emph{At $\gamma = 1$, the matrix $R_{1/2}^{-1}\,R_{1/2}^{\circ 2}$ is
upper triangular with eigenvalues $\{1, 3, 5, \ldots, 2N-1\}$.
In particular, $\kappa(1,N) = 1$ for all $N \geq 1$.}

\medskip\noindent
\textbf{Theorem~D} (Cross-sign vanishing and uniqueness of $\gamma=1$; \S\ref{sec:cross-sign}).
\emph{The cross-sign kernel satisfies
\begin{equation}\label{eq:intro-D}
  \Psi_\gamma(n,-m) = \tfrac{1}{2}\bigl(|n|^\gamma + |m|^\gamma
  - |n+m|^\gamma\bigr)
  \begin{cases}
    = 0 & \text{if } \gamma = 1, \\
    > 0 & \text{if } \gamma < 1, \\
    < 0 & \text{if } \gamma > 1,
  \end{cases}
\end{equation}
for all $n,m > 0$. Consequently, $\gamma = 1$ is the unique stability
index for which $\kappa_{\mathrm{real}}(\gamma) \geq 1$ on all real
trigonometric polynomials.}

\medskip\noindent
\textbf{Theorem~E} (Curvature with confining potential; \S\ref{sec:potential}).
\emph{For $L = -(-\Delta)^{\gamma/2} - V'(x)\,\partial_x$ on $\T$ with
$V(x) = -\omega^2\cos x$, the curvature on the first mode
$f = \cos x$ is
\begin{equation}\label{eq:intro-E}
  \kappa(\gamma,\omega^2) = \min_{x \in \T}
  \frac{1 + \alpha^2\cos 2x + \tfrac{\omega^2}{2}(\cos x + \beta\cos 3x)}
       {1 + \alpha\cos 2x},
\end{equation}
where $\alpha(\gamma) = 1 - 2^{\gamma-1}$ and
$\beta(\gamma) = (2^{\gamma+1} - 3^\gamma - 1)/2$.
At $\gamma = 1$: $\alpha = \beta = 0$, giving $\kappa = 1 - \omega^2/2$.}


\subsection{Context and prior work}

This paper provides the first positive Bakry--\'Emery curvature
result for the fractional Laplacian.

Spener, Weber, and Zacher~\cite{SpenerWeberZacher2020} showed that
$(-\Delta)^{\gamma/2}$ on $\R^d$ fails the curvature-dimension inequality
$\mathrm{CD}(\kappa,N)$ for all $\kappa$ and all finite $N$. Their result
leaves open the case of compact domains. Our Theorem~C provides
$\mathrm{CD}(1,\infty)$ on the torus at $\gamma = 1$: the compactness
(spectral gap $\lambda_1 = 1$) is essential.

Gentil and Imbert~\cite{GentilImbert2008} prove exponential convergence
of $\Phi$-entropies for the L\'evy--Fokker--Planck equation with rate
$\gamma$ for $\gamma$-stable processes with Ornstein--Uhlenbeck drift,
bypassing the chain rule via entropy production. Their approach gives
existence of convergence but not explicit Bakry--\'Emery constants.
Our results are complementary: we work on the torus rather than $\R^d$
but obtain exact curvature formulas controlled by the fBM covariance.

Chafa\"i and Malrieu (unpublished, cited in~\cite{GentilImbert2008},
Remark~4) observe that $\Gamma_2 \geq 0$ for pure-jump generators
follows from a double-integral representation as a square. We go
much further: we compute the \emph{full sesquilinear kernel}, identify
it as the Hadamard square of the fBM covariance, and extract
quantitative curvature bounds.


\subsection{Organization}

Section~\ref{sec:prelim}: preliminaries on the fractional Laplacian,
carr\'e du champ, and fBM.
Section~\ref{sec:duality}: the stable--fBM correspondence (identity
$\Psi_\gamma = R_{\gamma/2}$) and the distinguished point $\gamma = 1$.
Section~\ref{sec:gamma2}: Hadamard square identity (Theorem~A).
Section~\ref{sec:curvature}: curvature on positive frequencies
(Theorems~B and~C) and the phase transition.
Section~\ref{sec:real}: cross-sign structure and real polynomials
(Theorem~D).
Section~\ref{sec:potential}: curvature with confining potential
(Theorem~E).
Section~\ref{sec:func-ineq}: functional inequality consequences.
Section~\ref{sec:discussion}: spectral invariants and open problems.


%======================================================================
\section{Preliminaries}\label{sec:prelim}
%======================================================================

\subsection{The fractional Laplacian on the torus}

Let $\T = \R/(2\pi\Z)$ with $\mu(dx) = dx/(2\pi)$. The fractional
Laplacian $L_\gamma = -(-\Delta)^{\gamma/2}$ with $\gamma \in (0,2)$
acts as $L_\gamma(e^{inx}) = -|n|^\gamma\,e^{inx}$ for $n \in \Z$.
The semigroup is $P_t^\gamma(e^{inx}) = e^{-t|n|^\gamma}\,e^{inx}$,
with spectral gap $\lambda_1(\gamma) = 1$.

The space of trigonometric polynomials
$\mathcal{T} = \mathrm{span}\{e^{inx} : n \in \Z\}$ is a core
for $L_\gamma$ on $L^2(\T,\mu)$: it is dense, invariant under $P_t^\gamma$,
and contained in the domain of every power of $L_\gamma$.
Our curvature computations are performed on finite-dimensional
subspaces of $\mathcal{T}$; the resulting constants are therefore
exact on these subspaces and give upper bounds for the curvature
on the full domain.

\begin{remark}[Role of compactness]\label{rem:compact}
Spener, Weber, and Zacher~\cite{SpenerWeberZacher2020} proved
$(-\Delta)^{\gamma/2}$ on $\R^d$ fails $\mathrm{CD}(\kappa,N)$
for all $\kappa$ and all finite $N$. The torus provides the
translation-invariant compact setting where the Fourier basis
diagonalizes $L_\gamma$, making the curvature computation tractable.
The spectral gap $\lambda_1 = 1$ supplies the compactness that
enables positive curvature; whether analogous results hold on
other compact domains (e.g.\ spheres) remains open.
\end{remark}

The generator on $\R$ has the integral representation
\begin{equation}\label{eq:generator}
  L_\gamma f(x)
  = \int_{\R \setminus \{0\}}
  \bigl[f(x+z) - f(x) - f'(x)\,z\,\mathbf{1}_{|z| \leq 1}\bigr]
  \,\nu_\gamma(dz),
\end{equation}
where $\nu_\gamma(dz) = c_\gamma\,|z|^{-1-\gamma}\,dz$ is the
L\'evy measure.
On Fourier modes, the characteristic exponent is recovered from
the L\'evy--Khintchine formula:
\begin{equation}\label{eq:LK}
  |\xi|^\gamma
  = \int_{\R \setminus \{0\}}
  (1 - \cos \xi z)\,\nu_\gamma(dz).
\end{equation}


\subsection{Carr\'e du champ}

\begin{definition}\label{def:cdc}
The sesquilinear carr\'e du champ of $L_\gamma$ is
\begin{equation}\label{eq:cdc}
  \widetilde{\Gamma}_\gamma(f,g)(x)
  := \frac{1}{2} \int_{\R \setminus \{0\}}
  \bigl[f(x+z) - f(x)\bigr]\,
  \overline{\bigl[g(x+z) - g(x)\bigr]}\;
  \nu_\gamma(dz).
\end{equation}
On Fourier modes:
\begin{equation}\label{eq:psi-kernel}
  \widetilde{\Gamma}_\gamma(e^{i\xi\cdot}, e^{i\eta\cdot})(x)
  = e^{i(\xi-\eta)x}\,\Psi_\gamma(\xi,\eta),
\end{equation}
where $\Psi_\gamma(\xi,\eta)
= \tfrac{1}{2}(|\xi|^\gamma + |\eta|^\gamma - |\xi-\eta|^\gamma)$.
Exponentials $e^{i\xi\,\cdot}$ lie in the extended Dirichlet space
of~$L_\gamma$ on $\T$ for all $\gamma \in (0,2)$; see
\cite{BoettcherSchillingWang2013}, \S4.7, for the general framework
of carr\'e du champ for L\'evy-type operators.
\end{definition}

We also need the \emph{bilinear} carr\'e du champ, defined by replacing
the conjugate in~\eqref{eq:cdc} with the plain difference. On Fourier modes:
\begin{equation}\label{eq:psi-bilinear}
  \Gamma_\gamma^{\mathrm{bil}}(e^{inx}, e^{imx})(x)
  = e^{i(n+m)x}\,\Psi_\gamma^B(n,m),
  \qquad
  \Psi_\gamma^B(n,m) = \tfrac{1}{2}\bigl(|n|^\gamma + |m|^\gamma
  - |n+m|^\gamma\bigr).
\end{equation}
Note: $\Psi_\gamma^B(n,m) = \Psi_\gamma(n,-m)$.

\begin{definition}\label{def:gamma2}
The iterated carr\'e du champ is
\begin{equation}\label{eq:gamma2-def}
  \widetilde{\Gamma}_{\gamma,2}(f,f)
  := \frac{1}{2}\Bigl[
    L_\gamma\bigl(\widetilde{\Gamma}_\gamma(f,f)\bigr)
    - 2\,\widetilde{\Gamma}_\gamma(f, L_\gamma f)
  \Bigr].
\end{equation}
The Bakry--\'Emery condition is
$\widetilde{\Gamma}_{\gamma,2}(f,f)(x) \geq \kappa\,
\widetilde{\Gamma}_\gamma(f,f)(x)$.
\end{definition}


\subsection{Fractional Brownian motion}

For $H \in (0,1)$, fractional Brownian motion (fBM) with Hurst
parameter $H$ is a centered Gaussian process $B^H$ with covariance
\begin{equation}\label{eq:fbm-cov}
  R_H(s,t)
  = \frac{1}{2}\bigl(|s|^{2H} + |t|^{2H} - |s-t|^{2H}\bigr).
\end{equation}
The kernel $R_H$ is positive definite for all $H \in (0,1)$
(\cite{Mishura2008}, Proposition~1.2.1).
At $H = 1/2$: $R_{1/2}(s,t) = \min(s,t)$ for $s,t \geq 0$
(standard Brownian motion; independent increments).
For $H > 1/2$: positively correlated increments (long-range dependence).
For $H < 1/2$: negatively correlated increments (anti-persistence).


%======================================================================
\section{The Stable--fBM Correspondence}\label{sec:duality}
%======================================================================

\begin{theorem}[Stable--fBM correspondence]\label{thm:duality}
For $\gamma \in (0,2)$ and all $\xi, \eta \in \R$:
\begin{equation}\label{eq:duality}
  \Psi_\gamma(\xi,\eta)
  = \tfrac{1}{2}\bigl(|\xi|^\gamma + |\eta|^\gamma
  - |\xi - \eta|^\gamma\bigr)
  = R_{\gamma/2}(|\xi|, |\eta|).
\end{equation}
The carr\'e du champ of the $\gamma$-stable generator, evaluated
on Fourier modes, is the covariance of fBM with $H = \gamma/2$.
\end{theorem}

\begin{proof}
Let $f(x) = e^{i\xi x}$, $g(x) = e^{i\eta x}$.
Then $f(x+z) - f(x) = e^{i\xi x}(e^{i\xi z} - 1)$ and similarly
for $g$. By the sesquilinear carr\'e du champ~\eqref{eq:cdc}:
\begin{align*}
  \widetilde\Gamma_\gamma(f,g)(x)
  &= \frac{e^{i(\xi-\eta)x}}{2}\int_{\R \setminus \{0\}}
     (e^{i\xi z} - 1)(e^{-i\eta z} - 1)\,\nu_\gamma(dz).
\end{align*}
Expanding and using $\int (e^{i\lambda z} - 1)\,\nu_\gamma(dz)
= -|\lambda|^\gamma$ (from~\eqref{eq:LK} and symmetry of $\nu_\gamma$):
\[
  \Psi_\gamma(\xi,\eta)
  = \tfrac{1}{2}\bigl[|\xi|^\gamma + |\eta|^\gamma
  - |\xi - \eta|^\gamma\bigr].
\]
Setting $H = \gamma/2$ and noting $|s|^{2H} = |s|^\gamma$,
this is $R_H(|\xi|, |\eta|)$.
\end{proof}

\begin{remark}[Scope of the identification]\label{rem:scope}
Theorem~\ref{thm:duality} is an identity of sesquilinear kernels
on Fourier modes, not a Hilbert space isomorphism. We do not claim
that $\Gamma_\gamma(f,g) = \langle f,g \rangle_{\mathcal{H}}$ for
a Gaussian Hilbert space~$\mathcal{H}$ and general $f,g$.
The identification operates at the level of the kernel
$\Psi_\gamma(\xi,\eta)$; its power comes from the fact that the
Bakry--\'Emery curvature theory (Sections~\ref{sec:gamma2}--\ref{sec:real})
is entirely controlled by this kernel.
\end{remark}

\begin{remark}[Three literatures, one formula]\label{rem:three}
The kernel $\frac{1}{2}(|\xi|^\gamma + |\eta|^\gamma
- |\xi-\eta|^\gamma)$ appears in:
(i)~Schoenberg's theory of conditionally negative definite
functions~\cite{Schoenberg1938};
(ii)~Dirichlet forms for jump
processes~\cite{BoettcherSchillingWang2013};
(iii)~fBM covariance theory~\cite{Mishura2008}.
These three appearances have not previously been connected explicitly.
\end{remark}


\subsection{The distinguished point $\gamma = 1$}

\begin{proposition}\label{prop:self-dual}
At $\gamma = 1$ ($H = 1/2$):
$\Psi_1(\xi,\eta)
= \min(|\xi|, |\eta|)\,\mathbf{1}_{\xi\eta > 0}$.
This is the covariance of standard Brownian motion.
The Cauchy process ($\gamma = 1$) is thus the unique stable process
whose carr\'e du champ has the covariance structure of standard
Brownian motion.
For opposite-sign frequencies, $\Psi_1(\xi,-\eta) = 0$:
positive and negative frequencies decouple completely.
\end{proposition}

\begin{proof}
For $\xi, \eta > 0$:
$\frac{1}{2}(\xi + \eta - |\xi - \eta|) = \min(\xi, \eta)$.
For $\xi > 0, \eta < 0$:
$\frac{1}{2}(\xi - \eta - (\xi - \eta)) = 0$.
\end{proof}


\subsection{Mode coupling and the correspondence dictionary}

The normalized off-diagonal kernel
\[
  \rho_\gamma(\xi,\eta)
  = \frac{R_{\gamma/2}(|\xi|,|\eta|)}
    {|\xi|^{\gamma/2}\,|\eta|^{\gamma/2}}
\]
is the fBM increment correlation. It is increasing in $H = \gamma/2$:
for $\gamma > 1$ modes reinforce (long-range dependence);
for $\gamma < 1$ modes oppose (anti-persistence);
at $\gamma = 1$ modes are maximally decoupled.

The correspondence exchanges structurally opposite objects:
\[
\begin{array}{c|c}
\textbf{Stable process (index } \gamma\text{)}
& \textbf{fBM (Hurst } H = \gamma/2\text{)} \\[4pt]
\hline
\rule{0pt}{12pt}
\text{Markov, pure jump} & \text{non-Markov, continuous} \\
\text{Independent increments} & \text{Dependent increments} \\
\text{Characteristic exponent } |\xi|^\gamma &
  \text{Covariance } R_H(s,t) \\
\text{Carr\'e du champ kernel } \Psi_\gamma &
  \text{Covariance kernel } R_H
\end{array}
\]

The fBM interpretation is not merely a relabeling: it predicts and
explains the curvature phase transition at $\gamma = 1$.
Cross-sign vanishing (Theorem~\ref{thm:cross-sign}) is equivalent
to independence of fBM increments at $H = 1/2$; the block
diagonalization of the $\Psi_\gamma$ matrix
(Corollary~\ref{cor:block-diag}) is the Fourier manifestation
of this independence; and the curvature $\kappa = 1$
(Theorem~\ref{thm:cauchy}) follows because independent increments
produce the simplest possible covariance structure.


%======================================================================
\section{The $\Gamma_2$-Kernel Identity}\label{sec:gamma2}
%======================================================================

\begin{theorem}[Hadamard square]\label{thm:gamma2-kernel}
For $\gamma \in (0,2)$ and $\xi,\eta \in \R$:
\begin{equation}\label{eq:Phi-formula}
  \Phi_{\gamma,2}(\xi,\eta) = \bigl[\Psi_\gamma(\xi,\eta)\bigr]^2
  = \bigl[R_{\gamma/2}(|\xi|,|\eta|)\bigr]^2.
\end{equation}
\end{theorem}

\begin{proof}
For $f = \sum_{n \in S} a_n\,e^{i\xi_n x}$, expand
$\widetilde{\Gamma}_\gamma(f,f)(x)
= \sum_{n,m} a_n\,\bar{a}_m\,\Psi_\gamma(\xi_n,\xi_m)\,
  e^{i(\xi_n-\xi_m)x}$
and compute the two terms of~\eqref{eq:gamma2-def}.

\textbf{Term~1:}
$L_\gamma[\widetilde{\Gamma}_\gamma(f,f)] =
-\sum_{n,m} a_n\,\bar{a}_m\,|\xi_n-\xi_m|^\gamma\,
\Psi_\gamma(\xi_n,\xi_m)\,e^{i(\xi_n-\xi_m)x}$.

\textbf{Term~2:}
$\widetilde{\Gamma}_\gamma(f, L_\gamma f) =
-\sum_{n,m} a_n\,\bar{a}_m\,|\xi_m|^\gamma\,
\Psi_\gamma(\xi_n,\xi_m)\,e^{i(\xi_n-\xi_m)x}$.

Assembling and Hermitian-symmetrizing:
\[
  M_{nm}^{\mathrm{sym}}
  = \bigl[|\xi_n|^\gamma + |\xi_m|^\gamma - |\xi_n-\xi_m|^\gamma\bigr]\,
  \Psi_\gamma(\xi_n,\xi_m)
  = 2\,[\Psi_\gamma(\xi_n,\xi_m)]^2. \qedhere
\]
\end{proof}

\begin{remark}\label{rem:hadamard}
By the Schur product theorem, $\Psi_\gamma^{\circ 2}$ is positive
semidefinite whenever $\Psi_\gamma$ is (which holds for all
$\gamma \in (0,2)$ by the correspondence). This gives $\Gamma_{\gamma,2}(f,f) \geq 0$
as a special case of the observation of Chafa\"i--Malrieu
(cf.~\cite{GentilImbert2008}, Remark~4). The Hadamard square
identity~\eqref{eq:Phi-formula} is much stronger: it gives the
full kernel, not just the sign.
\end{remark}


%======================================================================
\section{Curvature on Positive Frequencies}\label{sec:curvature}
%======================================================================

\subsection{Matrix reduction}\label{sec:curv-pos}

\begin{theorem}\label{thm:curvature}
On the space $\mathcal{T}_N^+$ of positive-frequency trigonometric
polynomials $f = \sum_{n=1}^N a_n\,e^{inx}$, the Bakry--\'Emery
curvature is exactly
\begin{equation}\label{eq:kappa-def}
  \kappa(\gamma,N) = \lambda_{\min}\!\bigl(R_H^{\circ 2}, R_H\bigr),
  \qquad H = \gamma/2,
\end{equation}
where $R_H$ is the $N \times N$ fBM covariance matrix on $\{1,\ldots,N\}$.
Since $\mathcal{T} = \bigcup_N \mathcal{T}_N$ is a core for $L_\gamma$,
the curvature on the full domain satisfies
$\kappa(\gamma) \leq \lim_{N\to\infty} \kappa(\gamma,N)$.
\end{theorem}

\begin{proof}
Write $\mathbf{b}(x) = D(x)\,\mathbf{a}$ with
$D(x) = \mathrm{diag}(e^{ix},\ldots,e^{iNx})$. Then
$\widetilde{\Gamma}_\gamma(f,f)(x) = \mathbf{b}(x)^* R_H\,\mathbf{b}(x)$
and $\widetilde{\Gamma}_{\gamma,2}(f,f)(x) =
\mathbf{b}(x)^* R_H^{\circ 2}\,\mathbf{b}(x)$.
Since $\{D(x)\mathbf{a} : x \in \T, \mathbf{a} \in \C^N\} = \C^N$,
the infimum over $(x,\mathbf{a})$ equals
$\lambda_{\min}(R_H^{\circ 2}, R_H)$.
\end{proof}


\subsection{The Cauchy point $\gamma = 1$}\label{sec:cauchy}

\begin{theorem}\label{thm:cauchy}
At $\gamma = 1$ ($H = 1/2$), the matrix $R_{1/2}^{-1}\,R_{1/2}^{\circ 2}$
is upper triangular with diagonal entries $\{1, 3, 5, \ldots, 2N-1\}$.
Hence $\kappa(1,N) = 1$ for all $N$.
\end{theorem}

\begin{proof}
At $H=1/2$, $R_{1/2}(n,m) = \min(n,m)$. The inverse is tridiagonal
(discrete Laplacian). Direct computation: for $k < m$,
\[
  (R_{1/2}^{-1}\,R_{1/2}^{\circ 2})_{km}
  = -(k-1)^2 + 2k^2 - (k+1)^2 = -2,
\]
and $(R_{1/2}^{-1}\,R_{1/2}^{\circ 2})_{kk} = 2k-1$.
The matrix is upper triangular; eigenvalues are diagonal entries.
\end{proof}


\subsection{Phase transition}

\begin{theorem}\label{thm:phase-transition}
\begin{enumerate}[label=(\roman*)]
\item For $\gamma \in (0,1)$: $\kappa(\gamma) := \lim_N \kappa(\gamma,N)$
  exists, is finite, and satisfies $\kappa(\gamma) > 1$.
\item At $\gamma = 1$: $\kappa(1,N) = 1$ for all $N$.
\item For $\gamma \in (1,2)$: $\kappa(\gamma,N) \to -\infty$ as
  $N \to \infty$.
\end{enumerate}
\end{theorem}

Curvature values on positive frequencies (computed at $N=500$):

\medskip
\begin{center}
\begin{tabular}{ccc}
\toprule
$\gamma$ & $H$ & $\kappa_{\mathrm{pos}}(\gamma)$ \\
\midrule
$0.2$ & $0.10$ & $1.143$ \\
$0.4$ & $0.20$ & $1.164$ \\
$0.5$ & $0.25$ & $1.159$ \\
$0.8$ & $0.40$ & $1.100$ \\
$1.0$ & $0.50$ & $1.000$ \\
$1.2$ & $0.60$ & $0.871$ \\
$1.5$ & $0.75$ & $-0.19$ \\
\bottomrule
\end{tabular}
\end{center}

\begin{remark}\label{rem:pos-freq-only}
These values hold on positive-frequency polynomials only.
The curvature on real-valued trigonometric polynomials differs
and is treated in the next section.
\end{remark}


%======================================================================
\section{Real Polynomials and Cross-Sign Structure}\label{sec:real}
%======================================================================

For real-valued $f(x) = \sum_{n=1}^N [c_n\cos(nx) + s_n\sin(nx)]$,
the Fourier coefficients satisfy $a_{-n} = \bar{a}_n$, and the relevant
frequency set is $\{-N,\ldots,-1,1,\ldots,N\}$. The $\Psi_\gamma$ matrix
on this set has $2 \times 2$ block structure:
\begin{equation}\label{eq:block}
  \Psi_\gamma\Big|_{\{-N,\ldots,N\} \setminus \{0\}}
  = \begin{pmatrix} A & B \\ B^* & A \end{pmatrix},
\end{equation}
where $A_{nm} = \Psi_\gamma(n,m)$ for $n,m > 0$ (the positive-frequency
block) and $B_{nm} = \Psi_\gamma(n,-m)$ for $n,m > 0$ (the cross-sign
block).


\subsection{Cross-sign vanishing at $\gamma = 1$}\label{sec:cross-sign}

\begin{theorem}[Cross-sign vanishing]\label{thm:cross-sign}
For $n,m > 0$:
\begin{equation}\label{eq:cross-sign}
  \Psi_\gamma(n,-m) = \tfrac{1}{2}\bigl(n^\gamma + m^\gamma
  - (n+m)^\gamma\bigr).
\end{equation}
This vanishes identically if and only if $\gamma = 1$:
\begin{enumerate}[label=(\alph*)]
\item At $\gamma = 1$: $\Psi_1(n,-m)
  = \tfrac{1}{2}(n + m - (n+m)) = 0$ for all $n,m > 0$.
\item For $\gamma < 1$: $\Psi_\gamma(n,-m) > 0$
  (subadditivity of $|\cdot|^\gamma$).
\item For $\gamma > 1$: $\Psi_\gamma(n,-m) < 0$
  (superadditivity of $|\cdot|^\gamma$).
\end{enumerate}
\end{theorem}

\begin{proof}
Part~(a) is immediate from linearity of $|\cdot|^1$.
Parts~(b) and~(c): for $t \mapsto t^\gamma$ on $(0,\infty)$,
strict concavity ($\gamma < 1$) gives
$(n+m)^\gamma < n^\gamma + m^\gamma$, and strict convexity
($\gamma > 1$) gives the reverse inequality.
\end{proof}

\begin{corollary}\label{cor:block-diag}
At $\gamma = 1$, the cross-sign block $B = 0$, so the $\Psi_1$
matrix is block-diagonal. Positive and negative frequencies
decouple completely. For real $f$ with $a_{-n} = \bar{a}_n$,
$\kappa_{\mathrm{real}}(1,N) = 1$ for all $N$.
\end{corollary}


\subsection{Single-mode curvature formula}

\begin{proposition}\label{prop:single-mode}
For $f(x) = \cos(nx + \varphi)$:
\begin{equation}\label{eq:single-mode}
  \frac{\widetilde{\Gamma}_{\gamma,2}(f,f)(x)}
       {\widetilde{\Gamma}_\gamma(f,f)(x)}
  = |n|^\gamma \cdot
  \frac{1 + \alpha_n^2\,\cos(2nx+\theta)}
       {1 + \alpha_n\,\cos(2nx+\theta)},
\end{equation}
where $\alpha_n = 1 - 2^{\gamma-1}$ (independent of $n$)
and $\theta = 2\varphi$. The curvature is
\begin{equation}\label{eq:kappa1}
  \kappa_1(\gamma) = \frac{1 + \alpha^2}{1 + |\alpha|},
  \qquad \alpha = 1 - 2^{\gamma-1}.
\end{equation}
At $\gamma = 1$: $\alpha = 0$, giving $\kappa_1(1) = 1$.
For all $\gamma \neq 1$: $\kappa_1(\gamma) < 1$.
\end{proposition}

\begin{proof}
For $f = \cos(nx) = (e^{inx} + e^{-inx})/2$, the carr\'e du champ is
$\Gamma(f,f)(x) = \tfrac{1}{2}[1 + \alpha\cos(2nx)]$
with $\alpha = 1 - 2^{\gamma-1}$,
using $\Psi_\gamma(n,n) = n^\gamma$ and
$\Psi_\gamma(n,-n) = n^\gamma(1 - 2^{\gamma-1})$.
The $\Gamma_2$ computation yields the stated ratio.
The infimum over $x$ gives~\eqref{eq:kappa1}.
\end{proof}


\subsection{Uniqueness of $\gamma = 1$}

\begin{theorem}[Uniqueness on the torus]\label{thm:uniqueness}
On $\T$, $\gamma = 1$ is the unique stability index for which
$\Gamma_2 \geq \kappa\,\Gamma$ holds with $\kappa \geq 1$
on all real trigonometric polynomials. The three regimes are:
\begin{enumerate}[label=(\roman*)]
\item $\gamma < 1$: positive cross-sign coupling ($B > 0$),
  anti-persistent fBM dual ($H < 1/2$), $\kappa_{\mathrm{real}} < 1$.
\item $\gamma = 1$: zero cross-sign ($B = 0$),
  Brownian dual ($H = 1/2$), $\kappa_{\mathrm{real}} = 1$.
\item $\gamma > 1$: negative cross-sign ($B < 0$),
  long-range dependent dual ($H > 1/2$), $\kappa_{\mathrm{real}} < 1$.
\end{enumerate}
\end{theorem}

\begin{proof}
At $\gamma = 1$: Corollary~\ref{cor:block-diag} gives
$\kappa_{\mathrm{real}} = \kappa_{\mathrm{pos}} = 1$.
For $\gamma \neq 1$: Proposition~\ref{prop:single-mode} gives
$\kappa_1(\gamma) < 1$ already on the single mode $f = \cos x$.
\end{proof}

\begin{remark}[Probabilistic interpretation]
The cross-sign vanishing reflects the independence of fBM increments
at $H = 1/2$. In the Fourier setting, this independence manifests as
the orthogonality of positive and negative frequency contributions
to the carr\'e du champ.
\end{remark}


%======================================================================
\section{Curvature with Confining Potential}\label{sec:potential}
%======================================================================

Consider $L = -(-\Delta)^{\gamma/2} - V'(x)\,\partial_x$ on $\T$
with $V(x) = -\omega^2\cos x$, so that
$V'(x) = \omega^2\sin x$. The carr\'e du champ is unchanged;
the iterated carr\'e du champ acquires a drift correction:
\begin{equation}\label{eq:G2-drift}
  \Gamma_{L,2}(f,f) = \Gamma_{\gamma,2}(f,f)
  + \tfrac{1}{2}\bigl[
    b \cdot \nabla\Gamma(f,f) - 2\Gamma(f, b \cdot \nabla f)
  \bigr],
\end{equation}
where $b(x) = -V'(x) = -\omega^2\sin x$.

\begin{theorem}\label{thm:potential}
For $f(x) = \cos x$ and $V(x) = -\omega^2\cos x$:
\begin{equation}\label{eq:G2-ratio-potential}
  \frac{\Gamma_{L,2}(f,f)(x)}{\Gamma(f,f)(x)}
  = \frac{1 + \alpha^2\cos 2x + \tfrac{\omega^2}{2}(\cos x + \beta\cos 3x)}
         {1 + \alpha\cos 2x},
\end{equation}
where
$\alpha(\gamma) = 1 - 2^{\gamma-1}$ and
$\beta(\gamma) = (2^{\gamma+1} - 3^\gamma - 1)/2$.
\end{theorem}

\begin{proof}
The drift $b(x) = -\omega^2\sin x$ maps mode $n$ to modes $n \pm 1$.
The free part gives
$\Gamma_{\gamma,2}(f,f) = \frac{1}{2}[1 + \alpha^2\cos 2x]$
from Proposition~\ref{prop:single-mode}.
The drift correction assembles to
$\frac{\omega^2}{4}[\cos x + \beta\cos 3x]$,
using the key cancellation $\alpha + q = 1$ (where $q = 2^{\gamma-1}$),
which holds independently of $\gamma$.
\end{proof}

\begin{remark}
At $\gamma = 1$: $\alpha = \beta = 0$, so
$\Gamma_{L,2}/\Gamma = 1 + \frac{\omega^2}{2}\cos x$,
with minimum $1 - \omega^2/2$ at $x = \pi$.
\end{remark}

\begin{corollary}[Comparison with diffusion]\label{cor:comparison}
For the standard Laplacian ($\gamma = 2$) with the same potential:
$\kappa_{\mathrm{diff}} = -\omega^2$.
The jump curvature at $\gamma = 1$ is
$\kappa_{\mathrm{jump}} = 1 - \omega^2/2 > -\omega^2$
for all $\omega^2 > 0$. The jump process has better curvature because
the carr\'e du champ averages over the potential landscape.
\end{corollary}

The critical potential strength $\omega^2_{\mathrm{crit}}(\gamma)$
varies non-monotonically with $\gamma$:

\medskip
\begin{center}
\begin{tabular}{cc|cc}
\toprule
$\gamma$ & $\omega^2_{\mathrm{crit}}$ &
$\gamma$ & $\omega^2_{\mathrm{crit}}$ \\
\midrule
$0.3$ & $2.22$ & $1.0$ & $2.00$ \\
$0.5$ & $2.07$ & $1.1$ & $2.07$ \\
$0.7$ & $1.98$ & $1.2$ & $2.20$ \\
$0.8$ & $1.96$ & $1.3$ & $2.37$ \\
$0.9$ & $1.97$ & $1.5$ & $2.21$ \\
\bottomrule
\end{tabular}
\end{center}

\medskip\noindent
The non-monotonic behavior demonstrates that the fBM covariance
parameters $\alpha(\gamma)$ and $\beta(\gamma)$ genuinely control
functional inequality constants in the presence of confinement.

\begin{remark}[Generality of the potential]\label{rem:potential-general}
The choice $V(x) = -\omega^2\cos x$ is analytically convenient:
the drift $b(x) = \omega^2\sin x$ couples mode~$n$ only to
$n \pm 1$, yielding a closed-form curvature ratio.
For a general potential $V(x) = \sum_k v_k\,e^{ikx}$, the drift
couples mode~$n$ to all modes $n \pm k$, and the curvature
computation becomes a band-matrix eigenvalue problem of the same
type but with wider bandwidth. The method does not collapse;
it produces a larger but still computable linear algebra problem.
The cosine case captures the essential mechanism: the interplay
between the fBM parameters $\alpha(\gamma)$, $\beta(\gamma)$ and
the potential strength $\omega^2$.
\end{remark}


%======================================================================
\section{Functional Inequalities}\label{sec:func-ineq}
%======================================================================

On the flat torus ($V = 0$), the log-Sobolev constant
$C_{\mathrm{LS}} = 2$ for all $\gamma$, because the worst case
is concentrated on the first mode where
$\mathcal{E}_\gamma = 1$ independent of $\gamma$.

With a confining potential, Theorem~\ref{thm:potential} gives
curvature $\kappa(\gamma,\omega^2) > 0$ whenever
$\omega^2 < \omega^2_{\mathrm{crit}}(\gamma)$.
In this regime, the Poincar\'e inequality holds with constant
$1/\kappa(\gamma,\omega^2)$, and
$\mathcal{E}(P_t f) \leq e^{-2\kappa t}\,\mathcal{E}(f)$.

The $\gamma$-dependence means that different L\'evy indices give
genuinely different convergence rates for a fixed potential.
This cannot be obtained from the
Gentil--Imbert approach~\cite{GentilImbert2008}, which gives
a uniform rate independent of the potential shape.

The Bakry--\'Emery condition also implies the gradient estimate
\begin{equation}\label{eq:grad-est}
  \Gamma(P_t f, P_t f)(x) \leq e^{-2\kappa t}\,P_t[\Gamma(f,f)](x).
\end{equation}


%======================================================================
\section{Discussion and Open Problems}\label{sec:discussion}
%======================================================================

\subsection{Refined contraction and the spectral determinant at $\gamma = 1$}

The uniform gradient estimate~\eqref{eq:grad-est} uses only
the minimum eigenvalue $\kappa = 1$ of the curvature operator.
At $\gamma = 1$, the full eigenvalue spectrum
$\{\kappa_k = 2k-1\}_{k=1}^N$ (Theorem~\ref{thm:cauchy})
yields a sharper mode-by-mode estimate.

\begin{proposition}[Refined contraction]\label{prop:refined}
At $\gamma = 1$, decompose $f \in \mathcal{T}_N^+$ along the
eigenbasis $\{\phi_k\}_{k=1}^N$ of $M(1) = R_{1/2}^{-1}\,R_{1/2}^{\circ 2}$.
Then the $k$-th component of the carr\'e du champ contracts at rate
\begin{equation}\label{eq:refined-contraction}
  \Gamma(P_t f, P_t f)_k \leq e^{-2(2k-1)t}\,\Gamma(f,f)_k,
  \qquad k = 1,\ldots,N.
\end{equation}
The total energy satisfies
\begin{equation}\label{eq:total-energy}
  \int_\T \Gamma(P_t f, P_t f)\,d\mu
  \leq \sum_{k=1}^N e^{-2(2k-1)t}\,\mathcal{E}_k(f),
\end{equation}
where $\mathcal{E}_k(f)$ is the energy in the $k$-th eigendirection.
The uniform bound $e^{-2t}\,\mathcal{E}(f)$ is sharp only for
functions concentrated on the first mode ($k=1$).
\end{proposition}

The geometric mean contraction rate is
\[
  \Bigl(\prod_{k=1}^N \kappa_k\Bigr)^{1/N}
  = \bigl((2N-1)!!\bigr)^{1/N} \sim \frac{2N}{e}
  \qquad (N \to \infty),
\]
by Stirling's approximation.
For functions spread across $N$ modes, the effective decay rate
grows linearly in $N$, far exceeding the worst-case $\kappa = 1$.

\begin{corollary}[Volume comparison]\label{cor:volume}
The $\Gamma_2$-ellipsoid
$E_2 = \{f : \Gamma_2(f,f) \leq 1\}$ and the $\Gamma$-ellipsoid
$E_1 = \{f : \Gamma(f,f) \leq 1\}$ on $\mathcal{T}_N^+$ satisfy
\begin{equation}\label{eq:volume-ratio}
  \frac{\mathrm{vol}(E_1)}{\mathrm{vol}(E_2)}
  = (\det M)^{1/2} = \bigl((2N-1)!!\bigr)^{1/2}
  = \bigl(\E[X^{2N}]\bigr)^{1/2},
\end{equation}
where $X \sim \mathcal{N}(0,1)$.
At $\gamma = 1$, the $\Gamma_2$-form is $(2N-1)!!^{1/2}$ times
more concentrating than the $\Gamma$-form in $N$-dimensional volume.
\end{corollary}

\begin{proof}
For ellipsoids $\{v : v^T A v \leq 1\}$ and $\{v : v^T B v \leq 1\}$,
the volume ratio is $(\det A / \det B)^{-1/2} = (\det B^{-1}A)^{1/2}$.
Here $A = R_H$, $B = R_H^{\circ 2}$, and $\det M = \det(R_H^{-1}R_H^{\circ 2})
= \prod(2k-1) = (2N-1)!!$.
The identity $(2N-1)!! = \E[X^{2N}]$ for $X \sim \mathcal{N}(0,1)$
is standard.
\end{proof}

The unit determinant $\det R_{1/2} = 1$ (since $R_{1/2} = LL^T$
with $L$ unit lower triangular) characterizes $H = 1/2$ among
all fBM covariances, giving a further fingerprint of $\gamma = 1$.


\subsection{The Cauchy process as critical geometry}

The uniqueness of $\gamma = 1$ has four equivalent characterizations:
\begin{enumerate}[label=(\alph*)]
\item \emph{Curvature:} $\gamma = 1$ is the unique index with
  $\kappa_{\mathrm{real}} \geq 1$.
\item \emph{Algebraic:} $\Psi_1(n,-m) = 0$ (cross-sign vanishing),
  equivalent to linearity of $|\cdot|^1$.
\item \emph{Probabilistic:} the fBM dual has $H = 1/2$ (independent
  increments), so positive and negative frequencies decouple.
\item \emph{Determinantal:} $\det R_H = 1$ for all $N$; the
  $\Gamma$-kernel has unit determinant on every truncation.
\end{enumerate}


\subsection{Open problems}

\begin{enumerate}[label=(\roman*)]
\item \textbf{Curvature on $\R^d$.}
  Extend the torus computation to $\R^d$ with confining potentials.
  By~\cite{SpenerWeberZacher2020}, a confining potential is necessary
  for positive curvature.

\item \textbf{Sharp log-Sobolev constant.}
  With a confining potential, Theorem~\ref{thm:potential} shows
  $\gamma$-dependence of the curvature. Can the log-Sobolev constant
  for the full L\'evy--Fokker--Planck equation be expressed in terms of
  $\alpha(\gamma)$ and $\beta(\gamma)$?

\item \textbf{Closed-form curvature.}
  Is there a formula for
  $\kappa(\gamma) = \lim_N \lambda_{\min}(R_H^{\circ 2}, R_H)$
  on positive frequencies? Numerical values suggest a smooth
  function with maximum near $\gamma \approx 0.4$.

\item \textbf{General L\'evy processes.}
  For a general L\'evy process with measure $\nu$, the
  Fourier kernel
  $\Psi_\nu(\xi,\eta) = \frac{1}{2}\int (e^{i\xi z}-1)(e^{-i\eta z}-1)\,
  \nu(dz)$ is a positive definite kernel by Schoenberg's theorem.
  The stable case is special because scaling symmetry forces
  $\Psi_\nu$ to be self-similar, characterizing fBM among Gaussian
  processes~\cite{Mishura2008}.

\item \textbf{Multi-dimensional torus.}
  On $\T^d$, the direction-dependent fBM covariance should produce
  anisotropic curvature bounds.
\end{enumerate}


%======================================================================
\section*{Acknowledgments}
%======================================================================
AI tools were used as interactive assistants for mathematical
exploration, numerical computation, and drafting. The author
is responsible for all mathematical content and exposition.


%======================================================================
\begin{thebibliography}{20}

\bibitem{BakryGentilLedoux2014}
D.~Bakry, I.~Gentil, and M.~Ledoux.
\newblock \emph{Analysis and Geometry of Markov Diffusion Operators}.
\newblock Grundlehren Math. Wiss., vol.~348, Springer, 2014.

\bibitem{BoettcherSchillingWang2013}
B.~B\"ottcher, R.~Schilling, and J.~Wang.
\newblock \emph{L\'evy Matters III: L\'evy-Type Processes}.
\newblock Lecture Notes in Math., vol.~2099, Springer, 2013.

\bibitem{Chafai2004}
D.~Chafa\"{\i}.
\newblock Entropies, convexity, and functional inequalities.
\newblock \emph{J. Math. Kyoto Univ.}, 44(2):325--363, 2004.

\bibitem{Fontes2026a}
R.~Fontes.
\newblock Stochastic calculus as operator factorization.
\newblock \emph{arXiv preprint} arXiv:2601.09976 [math.PR], 2026.

\bibitem{Garofalo2019}
N.~Garofalo.
\newblock Fractional thoughts.
\newblock In \emph{New Developments in the Analysis of Nonlocal Operators},
  Contemp. Math., vol.~723, AMS, 2019.

\bibitem{GentilImbert2008}
I.~Gentil and C.~Imbert.
\newblock The {L}\'evy--{F}okker--{P}lanck equation: $\Phi$-entropies
  and convergence to equilibrium.
\newblock \emph{Asymptotic Analysis}, 59(3--4):125--138, 2008.

\bibitem{Mishura2008}
Y.~Mishura.
\newblock \emph{Stochastic Calculus for Fractional Brownian Motion
  and Related Processes}.
\newblock Lecture Notes in Math., vol.~1929, Springer, 2008.

\bibitem{Sato1999}
K.~Sato.
\newblock \emph{L\'evy Processes and Infinitely Divisible Distributions}.
\newblock Cambridge Univ. Press, 1999.

\bibitem{Schoenberg1938}
I.~J.~Schoenberg.
\newblock Metric spaces and positive definite functions.
\newblock \emph{Trans. Amer. Math. Soc.}, 44(3):522--536, 1938.

\bibitem{SpenerWeberZacher2020}
A.~Spener, F.~Weber, and R.~Zacher.
\newblock The fractional {L}aplacian has infinite dimension.
\newblock \emph{Comm. Partial Differential Equations}, 45(1):57--75, 2020.

\end{thebibliography}

\end{document}
